\section{研究方法}

% 写作提示：围绕研究目标定义估计量、响应变量、解释变量和统计假设，说明方法选择依据、模型形式、
% 参数估计、区间估计或检验流程及适用条件。下方统计方法与显著性水平仅为排版示例。

\subsection{描述性统计分析}
对数据进行描述性统计分析，包括均值、标准差、分位数等基本统计量。

\subsection{回归分析}
建立多元线性回归模型：
\begin{equation}
y = \beta_0 + \beta_1 x_1 + \beta_2 x_2 + \cdots + \beta_p x_p + \varepsilon
\end{equation}

\subsection{假设检验}
采用$t$检验和$F$检验对模型参数显著性进行检验（$\alpha = 0.05$）。
